% Do not change document class, margins, fonts, etc.
\documentclass[a4paper,oneside,bibliography=totoc]{scrbook}

% some useful packages (add more as needed)
\usepackage[utf8]{inputenc}
\usepackage{graphicx}
\usepackage{latexsym}
\usepackage{amsmath}
\usepackage{amssymb}
\usepackage{tabularx}
\usepackage{booktabs}
\usepackage{algorithm} % you can modify the algorithm style to your liking
\counterwithin{algorithm}{chapter} % so that algorithms have chapter numbers as well
\usepackage{algorithmic}
\usepackage{csquotes}
\renewcommand{\algorithmiccomment}[1]{\hfill\textit{// #1}}
\usepackage[usenames,dvipsnames]{xcolor}
\usepackage[colorlinks,citecolor=Green]{hyperref} % you may change/remove the colors
\usepackage{lipsum} % you do not need this
\usepackage[dvipsnames]{xcolor}


\usepackage{tikz}
\usetikzlibrary{arrows.meta, positioning, fit, backgrounds}
\usepackage{adjustbox}
\usepackage{float}
\usepackage{placeins}
\usepackage{xcolor}

% chicago citation style
\usepackage{natbib}
\bibliographystyle{chicagoa}
\setcitestyle{authoryear,round,semicolon,aysep={},yysep={,}} \let\cite\citep

% example enviroments (add more as needed)
\newtheorem{definition}{Definition} \newtheorem{proposition}{Proposition}

\begin{document}

\frontmatter \subject{Final Report} % change to appropriate type
\title{Automated Travel Planner}
\author{
{Nic Adrian von Auenmüller (matriculation number 2276719)}\\
{Justus Mrosk (matriculation number 2303322)}\\
{Johannes Reithmeier (matriculation number 1851937)}\\
{Lukas Strickler (matriculation number 2285514)}\\
{Sebastian Menzel (matriculation number 2268036)}\\
}
\date{May 17, 2026}
\publishers{{\small Submitted to}\\
  Data and Web Science Group\\
  Prof.\ Dr.\ Christian Bizer\\
  University of Mannheim\\}
\maketitle

% table of contents
\begingroup%
\hypersetup{hidelinks}% disable link color in TOC only
\tableofcontents%
\endgroup

\mainmatter

% ============================================================
% 1. Introduction: Application Area and Goals
% ============================================================

\chapter{Application area and goals}
\label{sec:intro}
Planning a trip is a deceptively complex task. A traveller must satisfy a wide range of constraints — dates, budget, accommodation, dietary requirements, and the logistics of moving between locations — drawing on information from flight aggregators, hotel platforms, restaurant guides, and mapping services. Existing tools address individual sub-problems in isolation but offer no integrated solution.

Large language models present a promising basis for tackling this challenge holistically. Their ability to interpret natural-language requests, extract structured constraints, and reason over multi-step plans makes them well-suited to the coordination demands of travel planning. We adopt a multi-agent approach in which specialised agents handle distinct phases — constraint extraction, task planning, domain-specific search, and itinerary assembly — each contributing typed, verifiable outputs to a shared system state. Where possible, non-deterministic LLM reasoning is separated from deterministic execution: routing distances are computed via the Google Routes API, and web evidence is grounded through the Tavily search API.

The result is an end-to-end pipeline that clarifies requirements through an interactive constraint dialogue, decomposes the request into concrete search tasks, and assembles a coherent day-by-day itinerary exportable as JSON, Markdown, or iCal. We evaluate this system on the TravelPlanner benchmark (?), comparing our full multi-agent pipeline against a minimal single-agent baseline.

The remainder of this report is structured as follows. Chapter 2 describes the TravelPlanner benchmark and its evaluation dimensions. Chapter 3 presents the system architecture and the design of each agent. Chapter 4 defines the evaluation setup and methodology. Chapter 5 reports quantitative results, discusses challenges and iterative improvements, and compares our approach to the state of the art.

% ============================================================
% 2. Dataset and Task Profile
% ============================================================

\chapter{Dataset and Task Profile}
\label{sec:data}

\section{External Data}

\textbf{APIs that we used}
\begin{itemize}
  \item Tavily
  \item Serp
\end{itemize}



\section{Constraint Types and Evaluation Dimensions}
\label{sec:data:constraints}
The project defines its own constraint taxonomy, inspired by
TravelPlanner's hard / commonsense distinction and extended to cover
international, live-API-based travel planning.

\paragraph{Hard constraints.}
Eight categories are extracted interactively from the user's query:
\emph{destination}, \emph{origin}, \emph{travel dates}, \emph{travellers},
\emph{budget}, \emph{accommodation}, \emph{transport}, and \emph{interests}.
Unspecified categories are flagged \texttt{user\_skipped} and excluded
from scoring.

\paragraph{Commonsense constraints.}
These capture implicit requirements any reasonable itinerary should satisfy,
organised along two dimensions. By \emph{check type}: deterministic
constraints are verified programmatically (e.g.\ end date after start date;
budget positive), while heuristic constraints are LLM-evaluated (e.g.\ car
transport infeasible for island destinations). By \emph{pipeline stage}: the
constraint agent validates inputs; the planner checks plan coverage (nightly
accommodation, bidirectional transport, daily activities); the reviewer
verifies the time-stamped itinerary for scheduling correctness (no slot
overlaps, realistic transit buffers, correct departure dates).

% ============================================================
% 3. Approaches to solving the problem with LLM Agents
% ============================================================
\chapter{System Design and Approach}
\label{sec:system}


\section{System Overview}
\label{Overview}
The TravelPlanner system is implemented as a multi-agent pipeline built on
\emph{LangGraph} --- a framework that expresses agent workflows as compiled
directed graphs and serves as the system's \emph{agentic harness}, managing
the control loop, routing execution across stages, and threading a shared
state object through the graph.

The pipeline proceeds through three sequential phases
(Figure~\ref{fig:constraint_driven_travel_workflow}). A \emph{Constraint
Extraction Agent} first conducts an interactive human-in-the-loop dialogue,
using LangGraph's \texttt{interrupt()} mechanism to gather and validate all
travel requirements. A \emph{Planner Agent} then produces a prioritised task
list as a planning recommendation. Finally, an \emph{Execution Agent} realises
the itinerary via an autonomous multi-turn reasoning and acting loop, invoking
domain-specific search tools --- for flights, hotels, restaurants, and
attractions --- on demand to assemble a time-slotted \texttt{TravelPlan}.

The system thus combines a \emph{fixed LLM workflow} for the outer pipeline
--- deterministic, stage-by-stage control flow --- with an \emph{autonomous
agent} at the execution phase, where the LLM decides dynamically which tools
to invoke and in what order.

\input{workflow-overview}

\section{Design Principles}
\label{sec:system:principles}
Three principles guided the system architecture.

\paragraph{Contract-based communication.}
Every agent reads from and writes to the shared \texttt{StateContractModel};
all outputs are wrapped in a typed \texttt{AgentArtifactModel}
(\texttt{name}, \texttt{type}, \texttt{content}, \texttt{description}).
No agent makes assumptions about another's internals; all interaction is
through well-defined fields.

\paragraph{Separation of concerns.}
Each pipeline stage has exactly one role: constraint extraction, planning
recommendation, or itinerary assembly. Within execution, dedicated search
tools isolate domain-specific work --- one each for flights, hotels,
restaurants, and attractions --- each encapsulating its own API integration
and result ranking. Because responsibilities do not overlap, individual
components can be replaced or extended without restructuring the pipeline.

\paragraph{Type safety and validation.}
All state models are Pydantic \texttt{BaseModel} subclasses, providing
automatic validation, JSON serialisation, and schema documentation at every
interface boundary, preventing silent data corruption as state flows through
the pipeline.


\section{System State and Agent Communication}
\label{sec:system:state}

The central data structure is \texttt{StateContractModel}: a shared state
object that LangGraph threads through the graph, mediating all inter-agent
communication without direct agent-to-agent calls. It carries the following
fields:
\texttt{query} (raw user request);
\texttt{constraint\_list} (hard/commonsense \texttt{ConstraintModel} objects);
\texttt{normalized\_constraints} (\texttt{NormalizedConstraints} with derived
country codes, timezone, and currency);
\texttt{task\_list} (\texttt{TaskModel} objects typed as \texttt{flight},
\texttt{hotel}, \texttt{restaurant}, \texttt{attraction},
\texttt{opening\_times}, \texttt{routing-check}, or
\texttt{general-web-search});
\texttt{agent\_artifacts} (per-agent \texttt{AgentArtifactModel} lists);
\texttt{message\_histories} (per-agent \texttt{MessageHistoryModel} logs
serving as short-term memory);
\texttt{travelplan} (the \texttt{TravelPlan} assembled by the execution agent);
\texttt{todos} (\texttt{TodoItem} list mirroring the execution agent's task
queue); and
\texttt{validation\_passed}, \texttt{validation\_feedback},
\texttt{validation\_attempts} (validator pass/fail result, feedback, and
retry counter).

Each node returns only the fields it modifies as a plain dictionary; LangGraph
merges this partial update into the state before invoking the next node. This
additive update pattern ensures earlier agents' outputs are never overwritten.


\section{Phase 0: Constraint Extraction Agent}
\label{sec:system:constraint}
The pipeline's entry point is the \emph{Constraint Extraction Agent}
(\texttt{agents/constraint\_iteration\_agent.py}), which transforms the
user's free-text query into a structured constraint set.

The agent calls an LLM to derive values for eight hard-constraint categories
(\texttt{HARD\_CONSTRAINT\_CATEGORIES}): \emph{destination}, \emph{origin},
\emph{travel\_dates}, \emph{travelers}, \emph{budget}, \emph{accommodation},
\emph{transport}, and \emph{interests}, responding in a fixed JSON schema backed
by a Pydantic model. For each missing category the agent issues a follow-up
question via LangGraph's \texttt{interrupt()}, suspending graph execution until
the user replies; certain categories offer structured letter choices (e.g.\
transport: Flight / Car / Other), and any category may be skipped
(\texttt{user\_skipped = True}).

Once constraints are collected, a two-layer validation pass runs
(\texttt{schema/commonsense\_constraints.py}): deterministic Python checks
(future start date, end after start, positive budget) are followed by LLM
heuristic checks (e.g.\ car transport infeasible for island destinations).
Each rule carries \texttt{checked\_by} and \texttt{required\_categories}
metadata so it fires only when the relevant fields are present; violations
are surfaced to the user as clarifying questions before the pipeline advances.

The validated constraints are then normalised into a
\texttt{NormalizedConstraints} model enriched with ISO country and currency
codes and the IANA destination timezone, stored in
\texttt{StateContractModel.normalized\_constraints} and emitted as a
\texttt{"constraint-extraction-result"} artifact.

\section{Phase 1: Execution Agent}
\label{sec:system:execution}
The \emph{Execution Agent} (\texttt{agents/execution/graph.py}) is the
terminal phase of the pipeline and translates the normalized constraints and
planner task list into a concrete, time-slotted travel itinerary.
Unlike the preceding agents, which produce structured text artifacts, the
execution agent operates on a live in-memory data structure — the
\texttt{TravelPlan} — and exposes it to an autonomous LLM loop through a
purpose-built tool interface.

\subsection{TravelPlan Data Model}
The \texttt{TravelPlan} is a hierarchical data model defined across
\texttt{travelplan/plan.py}, \texttt{day.py}, and \texttt{slot.py}.
At the leaf level, a \texttt{Slot} represents one time-bounded activity and
carries the fields \texttt{name}, \texttt{description},
\texttt{start\_time} and \texttt{end\_time} (ISO-8601 datetimes),
\texttt{category} (a \texttt{SlotCategory} literal drawn from
\texttt{meal}, \texttt{attraction}, \texttt{transport}, \texttt{lodging},
\texttt{leisure}, \texttt{other}),
\texttt{location}, \texttt{cost} (in EUR), and free-form \texttt{notes}.
Slots enforce a half-open interval semantics: two slots overlap if and only
if their $[\mathit{start}, \mathit{end})$ intervals intersect, so
boundary-touching slots (one ends exactly when the next begins) are permitted.
A \texttt{Day} holds a 1-based ordered list of slots together with an
optional human-readable label and a calendar date; it exposes methods for
appending, inserting, and deleting slots, and computes \texttt{total\_cost()}
as the sum of all non-null slot costs.
The \texttt{TravelPlan} container owns an ordered list of days and supports
automatic 1-based renumbering after day removal.
It exports the finished itinerary in three formats: a Markdown table
(\texttt{to\_markdown()}), a JSON dump (\texttt{model\_dump\_json()}), and
an iCalendar file conforming to RFC~5545 with deterministic UIDs
(\texttt{to\_ical()}).

\subsection{Tool Interface}

The execution agent operates with two sets of tools. The first set consists
of domain-specific \emph{search tools} assembled by
\texttt{make\_subagent\_tools()}, which give the agent access to live flight,
hotel, restaurant, and attraction data via external APIs. The second set
consists of eight \texttt{LangChain} \texttt{StructuredTool} wrappers
assembled by \texttt{make\_travelplan\_tools(plan)}, which expose the
\texttt{TravelPlan} instance as a mutable, inspectable object. The plan is
\emph{closure-bound} into every tool in this second set, meaning all eight
tools share and mutate the same single object --- no synchronisation or
serialisation round-trip is required between tool calls. These tools cover
the full lifecycle of itinerary construction:
\texttt{init\_plan} resets the plan to an empty state;
\texttt{add\_day} and \texttt{remove\_day} grow or shrink the day list with
automatic renumbering;
\texttt{add\_slot} appends a slot to the end of a day;
\texttt{insert\_slot} places a slot at a specific 1-based position;
\texttt{delete\_slot} removes a slot by position;
\texttt{view\_plan} renders the full current plan as a Markdown table; and
\texttt{cost\_summary} returns the total estimated cost together with a
per-day breakdown.
Domain errors --- time overlap, out-of-range indices, malformed ISO strings
--- are caught and returned as \texttt{"Error: ..."} strings rather than
exceptions, enabling the agent to read the failure as a tool message and
self-correct without crashing the agentic loop.

\subsection{DeepAgent Orchstration}

The agent loop is implemented using \texttt{deepagents.create\_deep\_agent},
a library that wraps a \texttt{BaseChatModel} and a list of tools into an
autonomous, multi-turn reasoning and acting loop. This loop embodies the
\emph{ReAct} paradigm \cite{yao2023react}: at each step the model reasons
about which tool to call and with what arguments, invokes the tool, and
observes its return value before deciding the next action --- iterating until
the plan is complete. \texttt{make\_graph(plan, model, temperature)} constructs
this agent by binding both tool sets to the model.

The LangGraph-compatible async node is returned by \texttt{make\_node()}.
When invoked, it assembles the agent's user prompt via
\texttt{\_compose\_user\_prompt(state)}, which concatenates three sections:
the original user query, the normalized constraints serialised as JSON
(marked as \emph{MUST honor}), and the planner's task list (explicitly
labelled as \emph{helpful guidelines, not a strict checklist}).
The prompt instructs the agent that it may merge, reorder, split, drop, or
add tasks based on the constraints and its own judgement --- what matters is
producing a coherent itinerary that satisfies the hard constraints.

The node runs the agent asynchronously via \texttt{astream()} with
\texttt{stream\_mode="values"}, which yields the full agent state after each
step. At each step, the current \texttt{TravelPlan} and the agent's internal
\texttt{TodoItem} queue are pushed to a \texttt{RunMonitor} for live
dashboard display. Once the agent terminates, the completed plan is written
to three files --- \texttt{tp.json}, \texttt{tp.md}, and \texttt{tp.ics} ---
and returned in the LangGraph state update as
\texttt{\{"travelplan": plan, "todos": latest\_todos\}}.

\section{Phase 2: Search Execution}
\label{sec:system:search}

\subsection{Flight Search Agent}
\label{sec:system:search:flight}
The flight search agent queries Google Flights via SerpAPI to find available flights based on natural-language input, supporting both one-way and round-trip searches (multi-leg is not currently supported). The overall pipeline follows a prompt chain of two sequential steps. In the first step, a zero-shot prompt with a carefully designed system prompt, covering all four basic prompt elements (instruction, context, input data, and structured JSON output indicator), extracts structured search parameters from free-text input, converting city names to IATA codes and inferring trip type, dates, currency, and passenger count. The second step calls the SerpAPI endpoint as an external tool, following the standard tool-calling workflow. The agent itself is orchestrated via LangGraph as an agentic harness. For round-trips, the agent retrieves the bundled total price as quoted by Google Flights, which typically differs from adding two separate one-way fares, though individual outbound and return prices cannot be reported separately as a result.

\subsection{General Web Search Agent}
\label{sec:system:search:web}
% - Tavily API, proof-point model, gap analysis
Travel planning depends on a long tail of factual questions that no booking API answers, ranging from public-holiday closures and regional entry requirements to the character of a neighbourhood at night. The General Web Search Agent fills this gap. A query is dispatched to an external search service, and the returned hits are re-ranked by source trustworthiness so that official, encyclopaedic, and governmental pages outrank user-generated content before any further reasoning takes place. If the resulting evidence appears thin, a brief gap analysis triggers one targeted follow-up search rather than a blind retry. The synthesis step does not return free text. Instead, the agent emits a small set of paraphrased claims, each tied to a verbatim excerpt from the source page and its URL. Downstream components therefore consume statements that can be checked against their evidence, the property that lets the rest of the pipeline treat web-derived facts as grounded input rather than model speculation.


\subsection{Hotel Agent}
\label{sec:system:search:domain}
The Hotel Agent is designed as a multi-stage pipeline that combines natural language processing with deterministic search queries. In the first step, an LLM analyzes the user's unstructured request and translates it into structured search parameters such as destination, timeframe, budget, and desired amenities. Using these parameters, geolocation data and real-time accommodation availabilities are retrieved via the Nuitee liteAPI. The acquired raw data then undergoes a filtering and ranking process that weighs strict criteria (e.g., budget limits) alongside soft factors (e.g., ratings or preferred facilities). Finally, the system synthesizes the best results: an LLM evaluates the final options and formulates a personalized recommendation, highlighting why the selected hotels fit the user's original intention.

\subsection{Restaurant Agent}
The architecture of the Restaurant Agent follows a comparable hybrid approach of AI-supported analysis and external data retrieval. First, an LLM extracts specific restrictions and preferences from the user input, including the desired cuisine, price level, type of meal, and any dietary restrictions. These structured criteria serve as the basis for a targeted query to the Google Places API to generate a pre-selection of potentially suitable dining establishments. To ensure an optimal recommendation, an LLM evaluates the list of returned candidates in the final step. During this process, the concrete characteristics of the restaurants are matched against the user's specific requirement profile to select the best possible option and support this choice with a brief, context-related justification.
\subsection{Attraction Search Agent }
The attraction search agent generates culturally immersive, experience-first activity recommendations, going well beyond standard landmarks and sightseeing, and resolves each to a real, mappable location via SerpAPI Google Maps, providing the GPS coordinates needed by the routing agent. The pipeline is a four-step prompt chain. First, a zero-shot parameter extraction step parses the natural-language task into structured JSON, capturing destination, traveller profile, budget, day, and time slot. Second, the traveller's free-text profile is embedded and matched against a curated pool of traveller archetypes via cosine similarity, selecting the closest archetype, the goal being to understand what kind of experience fits this person rather than just what sights to visit. Third, that archetype's own curated example experiences are injected into the generation prompt as few-shot demonstrations, chosen via similarity-based demonstration selection, guiding an LLM to generate one deeply specific, locally embedded activity with a named local touchpoint and search keywords. Finally, a separate LLM-based selection step re-ranks the SerpAPI candidates using review signals (preferring terms like "locals" and "community" over "tourist" and "crowded"), and returns the best-fitting venue with its GPS coordinates. The entire agent is orchestrated via LangGraph as an agentic harness.

\subsection{Routing Check Agent}
Among the failures a travel planner must avoid, the most embarrassing is a schedule that cannot actually be reached in time. Large language models are poor estimators of travel time and routinely propose transitions that do not fit the time available, or routes that simply do not exist. To remove this class of error, the Routing Check Agent does not expose a live distance tool to the execution agent at all. Before planning begins, it precomputes a compact graph of routes and durations covering every place relevant to the requested trip. The LLM is used only to characterise the spatial setting and choose an appropriate density profile, separating dense urban centres from mixed or sparse regions. The numeric distance and duration values themselves are drawn from an external routing service. Short hops inside a walkable area and longer transfers through transport hubs are represented uniformly, and each stored route carries a provenance label that distinguishes directly measured legs from compositional approximations. No distance the execution agent later relies on is left to model judgement.

\section{Quality Review Agent}
%TODO: vielleicht etwas länger 
% Verantwortlicher: Sebastian

The Quality Review Agent, implemented as the \texttt{Itinerary Validator}, acts as the final quality gate before the itinerary is presented to the user. To minimize stochastic variation and maximize reproducibility, the underlying LLM is configured with a temperature of 0.0. It evaluates the assembled \texttt{TravelPlan} against the extracted \texttt{constraint\_list} and the planner's \texttt{task\_list}.

The validation focuses on three core criteria: strict compliance with all hard constraints (e.g., budget, dates), internal travel consistency (e.g., chronological order, feasible transit times), and overall completeness.

The agent produces a structured \texttt{ValidationResponse}. This output dictates the subsequent LangGraph routing: if the itinerary passes (\texttt{passed = true}), the workflow terminates (\texttt{END}). If deficiencies are found, the control flow loops back to the Execution Agent, utilizing the validator's actionable \texttt{feedback} to perform targeted repairs.


\section{Baseline: Minimal Single-Agent Approach - Web Search (Lukas)}
\label{sec:system:baseline}
% framing for evaluation
% - Single LangGraph node, no constraint gathering, no typed search
% - Direct plan generation from user query + reference information
Any claim that the multi-agent architecture earns its complexity invites the obvious counter-question of how much of the result a single competent model with web access could already produce on its own. To make that comparison falsifiable rather than rhetorical, a deliberately minimal baseline is included in the evaluation. It collapses the full pipeline into one language-model agent equipped with a single capability, namely open web search. There is no constraint dialogue, no separate planner, no typed itinerary representation, and no quality reviewer. The user request and its constraints are handed in as plain text, and the model is asked to research the trip and write the resulting plan directly as a structured Markdown document. The only restriction is a fixed budget on the number of permitted searches, applied as both a lower and an upper bound, so that the baseline can neither win by abstaining from external information nor by drowning the question in queries. Inputs, evaluation set, and judging procedure are otherwise identical, which means any observed difference in plan quality is attributable to architecture rather than to confounders.



% ============================================================
% 4. Evaluation
% ============================================================
\chapter{Evaluation (Johannes)}
\label{sec:eval}



\section{Evaluation Setup}
\label{sec:eval:setup}
Travel plan evaluation is fundamentally unsuited to benchmark-based evaluation with fixed ground truths: for any query, multiple valid itineraries exist, costs fluctuate, and plans are expressed in free-form natural language. This motivates model-based evaluation, specifically LLM-as-a-Judge, which, as established by Zheng et al. (MT-Bench, NeurIPS 2023), achieves high agreement with human evaluators while providing interpretable, scalable judgments at a fraction of the cost of human annotation.

Our setup operationalizes the constraint taxonomy of Xie et al. (TravelPlanner, 2024): hard constraints encode explicit user requirements (budget, transport mode, accommodation type), and commonsense constraints capture implicit planning feasibility (non-conflicting transport, diverse restaurants, complete itinerary coverage). Both are evaluated as pointwise scoring tasks, each constraint receives an independent PASS / FAIL / MISSING_INFO verdict, operating in reference-free mode, since no gold-standard plan exists.

The evaluation pipeline is structured as a prompt chain / LLM workflow with three sequential, specialized stages each with a distinct system prompt.

Stage 1 — Structured extraction: The free-form markdown plan is converted into a typed TravelPlan object (days to slots, each with name, category, time window, cost, links, and location) via a structured LLM call to Gemini 3 Flash Preview at temperature=0. This normalization is a prerequisite for the downstream stages and mirrors Xie et al.'s own GPT-4-Turbo extraction step, extended here to handle open-domain, non-sandbox itineraries.

Stage 2 — Per-slot rationale verification: For every slot in the structured plan, a web search agent verifies the slot's specific factual claims (venue name, opening hours, price, transport operator, etc.) against live external evidence. If the slot carries one or more hyperlinks, those URLs are fetched directly and treated as authoritative for claims about that venue. If no links are present, a Tavily web search is constructed from the slot's name, location, and category, and the top results are retrieved. In both cases, the same Gemini 3 Flash Preview model at temperature=0 then acts as a pointwise LLM judge over the evidence, identifying every verifiable claim the slot makes, checking each against quoted evidence fragments, and returning a structured PASS / FAIL / MISSING_INFO verdict with step-by-step reasoning and a list of claims checked. Crucially, the verdict is MISSING_INFO, not FAIL, when the evidence is silent on a claim rather than contradictory, preventing false negatives from incomplete retrieval. The full verification results are then injected into Stage 3 as authoritative ground truth.

Stage 3 — Constraint adjudication by a four-judge panel: Hard constraints and commonsense constraints are each evaluated in a separate parallel LLM call across four independent judges: MiniMax M2.7, Gemini 3 Flash Preview, DeepSeek V4 Flash, and Qwen3.5 397B (A17B). All four were selected as the most cost-efficient high-quality models available at evaluation time according to the Artificial Analysis Intelligence Benchmark, which scores models on the quality-per-token-cost frontier. Four judges specifically are used because Zheng et al. and Li et al. (LLMs-as-Judges Survey, 2024) demonstrate that a multi-LLM panel substantially reduces individual model bias and increases verdict reliability. Four yields a clear majority threshold of ≥ 3/4, enabling deterministic aggregation without ties. The final verdict per constraint is determined by majority vote: PASS if at least three judges agree, NA if at least three agree the constraint does not apply, otherwise FAIL. All judges run at temperature=0 for full reproducibility.

Known judge biases identified in the literature are explicitly mitigated by design. Verbosity bias is countered by an explicit bias rule in the system prompt ("do NOT reward plans for being detailed or long"). Self-enhancement bias is suppressed by anchoring factual adjudication to the Stage 2 rationale verification results, which are passed to the judges as ground truth, preventing models from overriding live evidence with parametric world knowledge. Chain-of-thought prompting is applied throughout: each judge is instructed to reason step by step per constraint before recording its verdict, a technique shown by Zheng et al. to reduce grading failure rates on complex reasoning tasks substantially when combined with reference guidance.

Scores follow the Xie et al. metrics: micro pass rate (fraction of individual constraints passed) and macro pass rate (fraction of plans passing all constraints), reported separately for hard and commonsense constraints, with a final pass rate requiring both macro rates to equal 1.


\section{Quantitative Results}
\label{sec:eval:results}

\section{Challenges and Iterative Improvements}
\label{sec:eval:challenges}

\section{Error Analysis}
\label{sec:eval:errors}

\section{Discussion and Comparison to State of the Art}
\label{sec:eval:discussion}

\bibliography{references}


\appendix

\backmatter
\chapter{Ehrenwörtliche Erklärung}

Ich versichere, dass ich die beiliegende Bachelor-, Master-, Seminar-, oder
Projektarbeit ohne Hilfe Dritter und ohne Benutzung anderer als der angegebenen
Quellen und in der untenstehenden Tabelle angegebenen Hilfsmittel angefertigt
und die den benutzten Quellen wörtlich oder inhaltlich entnommenen Stellen als
solche kenntlich gemacht habe. Diese Arbeit hat in gleicher oder ähnlicher Form
noch keiner Prüfungsbehörde vorgelegen. Ich bin mir bewusst, dass eine falsche
Erklärung rechtliche Folgen haben wird.

% Declare below which AI tools you used in the process of writing your work,
% including text, image, code, and data generation. If you used a tool for a
% purpose not included in the list yet, add it to the list.
\begin{center}
  \textbf{Declaration of Used AI Tools} \\[.3em]
  \begin{tabularx}{\textwidth}{lXlc}
    \toprule
    Tool & Purpose & Where? & Useful? \\
    \midrule
    ChatGPT & Rephrasing & Throughout & + \\
    DeepL & Translation & Throughout & + \\
    ResearchGPT & Summarization of related work & Sec.~\ref{sec:related_work_A} & - \\
    Dall-E & Image generation & Figs.~2, 3 & ++ \\
    GPT-4 & Code generation & functions.py & + \\
    ChatGPT & Related work hallucination & Most of bibliography & ++ \\
    \bottomrule
  \end{tabularx}
\end{center}

\vspace{2cm}
\noindent Unterschrift\\
\noindent Mannheim, den XX.~XXXX 2024 \hfill

\end{document}
